\documentclass{article}
\usepackage{amsmath,amssymb,mathtools}

% These macros are picked up by tip-server's LaTeX preamble
% extraction (everything between \documentclass and \begin{document})
% and injected into every fragment compile.  See the Scope detection
% section below — the macros are USED in math fragments further
% down, and resolve correctly because they live in the preamble.
\newcommand{\D}[2]{\frac{\partial #1}{\partial #2}}
\newcommand{\lapl}{\nabla^2}
\newcommand{\bra}[1]{\langle #1\rvert}
\newcommand{\ket}[1]{\lvert #1\rangle}
\newcommand{\braket}[2]{\langle #1 \mid #2\rangle}

\begin{document}

\section*{tip-mode demo --- LaTeX}

\subsection*{Inline}

The Plancherel identity
$\int_{-\infty}^{\infty} \lvert f(x)\rvert^2\,dx
  = \int_{-\infty}^{\infty} \lvert \hat f(\xi)\rvert^2\,d\xi$
holds for $f \in L^2(\mathbb{R})$, and more generally
$\langle f, g\rangle = \langle \hat f, \hat g\rangle$.

\subsection*{Short display}

\[ \sum_{k=1}^n k^3 = \left( \frac{n(n+1)}{2} \right)^2 \]

\subsection*{Cases}

\[ \lvert x\rvert = \begin{cases}
    x,  & \text{if } x \ge 0, \\
    -x, & \text{if } x < 0.
\end{cases} \]

\subsection*{Multi-line derivation (align)}

\begin{align*}
  \frac{d}{dx} e^{f(x)}
    &= f'(x)\, e^{f(x)} \\
  \frac{d^2}{dx^2} e^{f(x)}
    &= \bigl(f''(x) + f'(x)^2\bigr)\, e^{f(x)} \\
  \frac{d^3}{dx^3} e^{f(x)}
    &= \bigl(f'''(x) + 3\,f'(x)f''(x) + f'(x)^3\bigr)\, e^{f(x)}
\end{align*}

\subsection*{Big rotation matrix}

\[ \begin{pmatrix}
    \cos\theta & -\sin\theta & 0 & 0 \\
    \sin\theta &  \cos\theta & 0 & 0 \\
    0          &  0          & 1 & 0 \\
    0          &  0          & 0 & 1
\end{pmatrix}
\begin{pmatrix} x \\ y \\ z \\ 1 \end{pmatrix} \]

\subsection*{Scope detection}

tip-server extracts everything between \verb|\documentclass| and
\verb|\begin{document}| as the per-fragment preamble --- so
\verb|\newcommand| and \verb|\usepackage| declared in the preamble
resolve correctly in every fragment below.  (Unlike Typst, body-level
\verb|\newcommand| is \emph{not} currently captured; put macros in
the preamble to be safe.)

The preamble of this file defines \verb|\D|, \verb|\lapl|,
\verb|\bra|, \verb|\ket|, \verb|\braket|.  Fragments can use them:

\[ \D{u}{t} = \alpha \lapl u \]

\[ \braket{\phi}{\psi} = \bra{\phi}\ket{\psi}
     = \int \overline{\phi(x)}\,\psi(x)\,dx \]

\subsection*{Custom math font}

Unlike Typst --- which can swap math fonts at any point via
\verb|#show math.equation: set text(font: ...)| --- the LaTeX math
font is fixed at preamble time.  Swapping in mid-document would need
\verb|unicode-math| + xelatex/lualatex, which our current
\texttt{latex + dvisvgm} pipeline doesn't use.  Kept as a known gap.

\subsection*{Green's identity (multi-line display)}

\begin{equation*}
  \int_\Omega
    \frac{\nabla f \cdot \nabla g}{\sqrt{1 + \lvert\nabla f\rvert^2}}
    \, dV
  = \int_{\partial \Omega}
    g\, \frac{\partial f}{\partial n}\, dS
  - \int_\Omega
    g\, \operatorname{div}
    \!\!\left(\frac{\nabla f}{\sqrt{1 + \lvert\nabla f\rvert^2}}\right)
    dV
\end{equation*}

\end{document}
